\documentclass[11pt,a4paper]{article}
\usepackage[T1]{fontenc}
\usepackage[utf8]{inputenc}
\usepackage[margin=2.2cm]{geometry}
\usepackage{amsmath,amssymb}
\usepackage{enumitem}
\usepackage{hyperref}
\usepackage{xcolor}
\hypersetup{colorlinks=true,linkcolor=blue,urlcolor=blue}

\setlength{\parskip}{0.35em}
\setlength{\parindent}{0pt}

\title{\vspace{-2.2em}\textbf{0--k BFS: Bounded Integer Weights Shortest Path Problem}\\
\large Project Plan -- Algorithms Course, C++ Implementation}
\author{Artemov Makar, Mark Shkut \\ Skoltech, Algorithms 2026}
\date{May 13, 2026}

\begin{document}
\maketitle
\vspace{-1.5em}

\section*{1. Motivation and Problem Statement}
The classical single-source shortest path (SSSP) problem on a weighted graph
$G=(V,E,w)$ with non-negative weights is solved by Dijkstra's algorithm in
$O((V+E)\log V)$ with a binary heap, or $O(E + V\log V)$ with a Fibonacci heap.
However, when edge weights are restricted to a small set of non-negative
integers $\{0,1,\dots,k\}$ -- a setting that arises frequently in grid path
planning, layered graphs, 2-SAT-like reductions, transport networks with a few
toll levels and many competitive-programming problems -- one can replace the
priority queue with a structure of $O(k)$ FIFO queues and obtain a strictly
faster $O(kV+E)$ algorithm. The aim of this project is to implement, analyse
and empirically compare the family of \emph{0--k BFS} algorithms against
Dijkstra-style baselines in C++.

\section*{2. Objectives}
\begin{itemize}[leftmargin=1.4em,itemsep=0.1em,topsep=0.1em]
    \item Implement a reusable C++17 graph library with three SSSP backends:
          (i) 0--1 BFS using \texttt{std::deque}; (ii) 0--k BFS via edge
          subdivision; (iii) 0--k BFS via multiple FIFO queues with the
          circular $(k{+}1)$-queue optimisation.
    \item Implement two Dijkstra baselines: \texttt{std::priority\_queue} (binary
          heap) and an indexed $d$-ary heap.
    \item Prove correctness and derive the running time / memory of each
          variant; identify exactly when 0--k BFS beats Dijkstra.
    \item Benchmark on synthetic and real graph families and report results.
\end{itemize}

\section*{3. Algorithms to Be Implemented}
\textbf{(a) 0--1 BFS.} Edges have weights in $\{0,1\}$. Replace the BFS queue
with a deque: relax an edge $(v,u)$ of weight $0$ by \texttt{push\_front}, of
weight $1$ by \texttt{push\_back}. The inf-check of plain BFS is replaced by
the relaxation $dist[v]+w(v,u)<dist[u]$, which is required because a vertex
may be reached several times before its final distance is fixed. Complexity
$O(V+E)$, memory $O(V)$.

\textbf{(b) 0--k BFS via edge subdivision.} Each edge of weight $w>1$ is split
into $w$ unit edges through $w-1$ fresh dummy vertices, reducing the problem
to 0--1 BFS on a graph with $V' \le (k-1)E+V$ vertices and $E' \le kE$ edges.
Total complexity $O(kE+V)$, memory $O(kE+V)$.

\textbf{(c) 0--k BFS via multiple queues.} Maintain queues $Q_0,Q_1,\dots$
where $Q_i$ holds vertices currently at $dist=i$, and process them in
non-decreasing order. A naive bound needs $k(V-1)$ queues; since at any moment
at most $k+1$ queues are non-empty (a relaxation moves a vertex by at most $k$
steps forward), it suffices to keep a circular buffer of $k+1$ queues indexed
by $dist[v] \bmod (k+1)$. Complexity $O(kV+E)$, memory $O(k+E)$.

\textbf{(d) Dijkstra baselines} with \texttt{std::priority\_queue} and a
hand-rolled 4-ary indexed heap supporting \texttt{decrease\_key}.

\section*{4. Theoretical Comparison}
\begin{center}
\begin{tabular}{lll}
\hline
Algorithm & Time & Memory \\ \hline
BFS (unit weights)              & $O(V+E)$           & $O(V)$ \\
0--1 BFS (deque)                & $O(V+E)$           & $O(V)$ \\
0--k BFS, subdivision           & $O(kE+V)$          & $O(kE+V)$ \\
0--k BFS, $k{+}1$ circular queues& $O(kV+E)$         & $O(k+E)$ \\
Dijkstra, binary heap           & $O((V+E)\log V)$   & $O(V+E)$ \\
Dijkstra, Fibonacci heap        & $O(E + V\log V)$   & $O(V+E)$ \\ \hline
\end{tabular}
\end{center}
For fixed $k$ and $E=\Theta(V)$ both 0--k variants are asymptotically
$\log V$ faster than a heap-based Dijkstra; the gap shrinks once $k = \omega(\log V)$.

\section*{5. Experimental Plan}
\begin{itemize}[leftmargin=1.4em,itemsep=0.1em,topsep=0.1em]
    \item \textbf{Graph families:} (i) random $G(n,p)$ with weights drawn from
          $\{0,\dots,k\}$; (ii) 4-connected grids with terrain costs (a classic
          0--k case); (iii) layered DAGs; (iv) DIMACS USA road network with
          weights quantised to $k\in\{4,8,16\}$; (v) adversarial chain graphs
          that force many relaxations.
    \item \textbf{Scale target.} Following the reviewer's note we aim for
          $V$ as close to $10^9$ as the workstation allows. Explicit-edge
          backends top out around $V \approx 10^8$ (a CSR with $E=4V$ and
          \texttt{uint32\_t} indices already costs $\approx 20$\,GB);
          \emph{implicit grid graphs} make $V=10^9$ feasible because the
          neighbour function is computed in $O(1)$ and only the distance
          array ($\approx 4$\,GB at \texttt{uint32\_t}) lives in memory.
    \item \textbf{Parameters:} $V \in \{10^6, 10^7, 10^8, 10^9\}$ for grids,
          $V \in \{10^6, 10^7\}$ for general graphs;
          $k \in \{1,2,4,8,16,64\}$; edge density sparse / dense.
    \item \textbf{Metrics:} wall-clock time (Google Benchmark),
          \texttt{perf}-counted instructions and L1/LLC misses, peak RSS,
          number of relaxations and queue operations.
    \item \textbf{Validation:} every backend's output is checked against the
          \texttt{std::priority\_queue} Dijkstra on every test instance.
\end{itemize}

\section*{6. Expected Limitations}
0--k BFS requires (a) non-negative integer weights with a small known bound
$k$, (b) a graph small enough that $O(kE)$ subdivision memory fits, or the
circular-queue trick is used. It does not generalise to real-valued weights,
to negative edges (Bellman--Ford / SPFA territory), or to problems requiring
the full SSSP shortest-path DAG with tie-breaking by hop count. For large
$k$ ($k \gtrsim \log V$) heap-based Dijkstra is the better choice.

\section*{References}
\small
[1] T.~H. Cormen, C.~E. Leiserson, R.~L. Rivest, C.~Stein. \emph{Introduction to Algorithms}, 4th ed., MIT Press, 2022. \quad
[2] E.~W. Dijkstra. ``A note on two problems in connexion with graphs.'' \emph{Numer. Math.} 1 (1959), 269--271. \quad
[3] R.~K. Ahuja, K.~Mehlhorn, J.~B. Orlin, R.~E. Tarjan. ``Faster algorithms for the shortest path problem.'' \emph{J.\ ACM} 37(2), 1990. \quad
[4] CP-Algorithms team. ``0--1 BFS,'' \texttt{cp-algorithms.com/graph/01\_bfs.html}.

\end{document}
